
\section{Problem 64: The Erd\H{o}s--Gy\'arf\'as conjecture (power-of-two cycle lengths)}

\subsection*{1) ROUND-4 OBJECTIVE}
\textbf{Path (A) --- proof, via a sharper minimal-counterexample reduction.} Round~3 reduced any putative counterexample to a connected edge-minimal graph $H$ with $\delta(H)=3$ and $V_{\ge 4}(H)$ independent, and obtained the quantitative bound that at least $4/7$ of the vertices have degree exactly $3$. The most promising next step is to strengthen that reduction itself: show that edge-minimality forces \emph{interaction among the degree-$3$ vertices}, not just a weak counting inequality. This succeeds. I prove that every degree-$3$ vertex of $H$ has a degree-$3$ neighbor, which upgrades the $4/7$ bound to a sharp $2/3$ bound and isolates the equality case.

\subsection*{2) Round-3 FOUNDATION USED}
I use the following Round~3 results without reproving them:
\begin{itemize}
  \item \textbf{Round-3 Lemma~5.1 and Corollary~5.2.} If a counterexample exists, then there is a connected edge-minimal counterexample $H$ such that $\delta(H)=3$, and the set
  \[
    B:=V_{\ge 4}(H)=\{v\in V(H): d_H(v)\ge 4\}
  \]
  is independent.
  \item \textbf{Round-3 Lemma~5.3.} With
  \[
    A:=V_3(H)=\{v\in V(H): d_H(v)=3\},
  \]
  one has $|A|\ge 4|V(H)|/7$.
  \item \textbf{Round-2/3 validated literature inputs.} Any counterexample has diameter at least $3$; every $P_{13}$-free graph with minimum degree at least $3$ satisfies the conjecture; hence every counterexample contains an induced $P_{13}$; any counterexample has at least $17$ vertices; and any bipartite counterexample has at least $32$ vertices.
\end{itemize}

\subsection*{3) NEW INSIGHT / TOOL (ROUND-4)}
The new tool is a \textbf{subgraph-minimality upgrade} of Round~3's edge-minimality. For the chosen edge-minimal counterexample $H$, \emph{no proper subgraph} of $H$ can still have minimum degree at least $3$, because every proper subgraph has fewer edges and cannot create new cycles.

This stronger minimality has three consequences.
\begin{enumerate}
  \item Every degree-$3$ vertex has a degree-$3$ neighbor.
  \item The induced subgraph $H[A]$ has minimum degree at least $1$.
  \item Consequently,
  \[
    |A|\ge 2|B|,
  \]
  so at least $2/3$ of all vertices have degree exactly $3$.
\end{enumerate}
Moreover, equality is rigid: if $|A|=2|B|$, then $H[A]$ is a perfect matching and every vertex of $B$ has degree exactly $4$.

\subsection*{4) ATTACK PLAN (ROUND-4)}
After Round~3, the exact unresolved gap was that the edge-minimal core still allowed a substantial population of vertices of degree at least $4$; the bound $|A|\ge 4|V(H)|/7$ was too weak to make the counterexample class genuinely narrow.

The Round~4 plan is therefore:
\begin{itemize}
  \item strengthen Round~3 Lemma~5.1 from \emph{single-edge deletions} to \emph{arbitrary proper subgraphs};
  \item use that stronger minimality to prove that vertices of degree $3$ cannot be supported entirely by higher-degree neighbors;
  \item convert this local structural fact into a global counting theorem $|A|\ge 2|B|$;
  \item isolate the equality case as a sharply formulated obstruction.
\end{itemize}
This overcomes the main Round~3 obstacle by forcing the degree-$3$ vertices to form a nontrivial internal backbone.

\subsection*{5) WORK (ROUND-4)}

\paragraph{Lemma 5.9 (new: subgraph-minimality of the edge-minimal counterexample).}
Let $H$ be the connected edge-minimal counterexample provided by Round~3 Lemma~5.1. Then no proper subgraph $F\subsetneq H$ satisfies $\delta(F)\ge 3$.

\emph{Proof.}
Suppose that some proper subgraph $F\subsetneq H$ has $\delta(F)\ge 3$. Since $F$ is a proper subgraph, it has strictly fewer edges than $H$. Also, every cycle of $F$ is a cycle of $H$, because subgraphs do not create new cycles. Since $H$ has no cycle whose length is a power of $2$, neither does $F$. Therefore $F$ is also a counterexample, contradicting the choice of $H$ as a counterexample with the minimum possible number of edges. \qed

\paragraph{Corollary 5.10 (new: every degree-$3$ vertex meets $A$).}
Every vertex of $A$ has at least one neighbor in $A$. Equivalently, the induced subgraph $H[A]$ has minimum degree at least $1$.

\emph{Proof.}
Let $a\in A$. Suppose, for contradiction, that every neighbor of $a$ lies in $B$. Since $d_H(a)=3$, the graph $H-a$ is a proper subgraph of $H$. Every neighbor of $a$ had degree at least $4$ in $H$, so after deleting $a$ each of those neighbors still has degree at least $3$. All other vertex degrees are unchanged. Hence
\[
  \delta(H-a)\ge 3.
\]
This contradicts Lemma~5.9. Therefore $a$ has a neighbor in $A$. Since $a\in A$ was arbitrary, $\delta(H[A])\ge 1$. \qed

\paragraph{Theorem 5.11 (new: the $2/3$ theorem for the edge-minimal core).}
With $A=V_3(H)$ and $B=V_{\ge 4}(H)$ as above,
\[
  |A|\ge 2|B|.
\]
Equivalently,
\[
  |A|\ge \frac{2}{3}|V(H)|,
  \qquad
  |B|\le \frac{1}{3}|V(H)|.
\]
In particular, at least two-thirds of the vertices of $H$ have degree exactly $3$.

\emph{Proof.}
By Corollary~5.10, every vertex of $H[A]$ has degree at least $1$ inside $H[A]$. Hence
\[
  2e_H(A)\ge |A|.
\]
Every vertex of $A$ has total degree $3$, so
\[
  3|A| = 2e_H(A)+e_H(A,B).
\]
Using $2e_H(A)\ge |A|$, we obtain
\[
  e_H(A,B) = 3|A|-2e_H(A) \le 2|A|.
\]
On the other hand, Round~3 Corollary~5.2 says that $B$ is independent. Therefore every edge incident with a vertex of $B$ lies in $E_H(A,B)$, and since each vertex of $B$ has degree at least $4$,
\[
  e_H(A,B)=\sum_{b\in B} d_H(b)\ge 4|B|.
\]
Combining the two inequalities yields
\[
  4|B|\le e_H(A,B)\le 2|A|,
\]
so $|A|\ge 2|B|$. Since $|V(H)|=|A|+|B|$, this is equivalent to
\[
  |A|\ge \frac{2}{3}|V(H)|,
  \qquad
  |B|\le \frac{1}{3}|V(H)|.
\]
\qed

\paragraph{Corollary 5.12 (new: rigidity of the equality case).}
Equality $|A|=2|B|$ in Theorem~5.11 holds if and only if both of the following are true:
\begin{enumerate}
  \item every vertex of $B$ has degree exactly $4$;
  \item the graph $H[A]$ is $1$-regular, i.e.
  $H[A]$ is a disjoint union of edges (a perfect matching on $A$).
\end{enumerate}

\emph{Proof.}
If $|A|=2|B|$, then the proof of Theorem~5.11 gives
\[
  4|B|\le e_H(A,B)\le 2|A|=4|B|,
\]
so equality holds throughout. Hence
\[
  e_H(A,B)=4|B|,
\]
which forces $d_H(b)=4$ for every $b\in B$, because each $b\in B$ has degree at least $4$ and
\[
  e_H(A,B)=\sum_{b\in B} d_H(b).
\]
Also,
\[
  e_H(A,B)=2|A|.
\]
Substituting into
\[
  3|A|=2e_H(A)+e_H(A,B)
\]
shows that
\[
  2e_H(A)=|A|.
\]
Since Corollary~5.10 says that every vertex of $H[A]$ has degree at least $1$, the equality $2e_H(A)=|A|$ implies that every vertex of $H[A]$ has degree exactly $1$. Thus $H[A]$ is $1$-regular, hence a perfect matching.

Conversely, if every vertex of $B$ has degree $4$ and $H[A]$ is a perfect matching, then
\[
  2e_H(A)=|A|,
  \qquad
  e_H(A,B)=3|A|-2e_H(A)=2|A|,
\]
while also
\[
  e_H(A,B)=\sum_{b\in B} d_H(b)=4|B|.
\]
Therefore $2|A|=4|B|$, i.e.
\[
  |A|=2|B|.
\]
\qed

\paragraph{Lemma 5.13 (new: no component of $H[A]$ is closed under $B$-neighbors).}
Assume $H[A]$ has at least two connected components. Then for every connected component $C$ of $H[A]$, there exists a vertex $b\in B$ that has a neighbor in $C$ and a neighbor in $A\setminus V(C)$.

\emph{Proof.}
Fix a connected component $C$ of $H[A]$, and let
\[
  S:=V(C)\cup \bigl(N_H(V(C))\cap B\bigr).
\]
Because $C$ is a connected component of $H[A]$, every neighbor in $A$ of a vertex of $C$ also lies in $C$. Thus every vertex of $C$ has all three of its neighbors inside $S$, so it has degree $3$ in $H[S]$.

Suppose, toward a contradiction, that every $b\in B\cap S$ has all its neighbors in $C$. Then each such $b$ keeps all of its neighbors in $H[S]$, and therefore has degree at least $4$ in $H[S]$. Hence
\[
  \delta(H[S])\ge 3.
\]
Since $H[A]$ has at least two connected components, there is a vertex of $A\setminus V(C)$, so $S\neq V(H)$. Therefore $H[S]$ is a proper induced subgraph of $H$, contradicting Lemma~5.9.

Thus some $b\in B\cap S$ must have a neighbor outside $C$. Since $B$ is independent, that outside neighbor must lie in $A\setminus V(C)$. \qed

\paragraph{Updated obstruction profile after Round~4.}
Combining the new Round~4 results with the validated inputs from Rounds~2 and~3, we obtain the following.

If a counterexample exists, then there is a connected edge-minimal counterexample $H$ such that:
\begin{enumerate}
  \item $\delta(H)=3$;
  \item $B=V_{\ge 4}(H)$ is independent;
  \item every vertex of $A=V_3(H)$ has a neighbor in $A$;
  \item at least $2/3$ of the vertices of $H$ lie in $A$;
  \item if equality $|A|=2|B|$ holds, then $H[A]$ is a perfect matching and every vertex of $B$ has degree exactly $4$;
  \item if $H[A]$ is disconnected, then every component of $H[A]$ is linked to another component through some vertex of $B$;
  \item $\operatorname{diam}(H)\ge 3$;
  \item $H$ contains an induced $P_{13}$;
  \item $|V(H)|\ge 17$;
  \item if $H$ is bipartite, then $|V(H)|\ge 32$.
\end{enumerate}
Thus the unresolved core is now narrower than in Round~3: a putative counterexample must be \emph{degree-$3$-dominated by a full two-thirds margin}, with the degree-$3$ vertices themselves forming a nontrivial internal scaffold.

\subsection*{6) ADVERSARIAL VERIFICATION}
\begin{itemize}
  \item \textbf{Does Lemma~5.9 really strengthen Round~3 Lemma~5.1?} Yes. Round~3 handled deletion of a single edge whose endpoints both had degree at least $4$. Lemma~5.9 applies to \emph{every} proper subgraph. The key monotonicity fact is valid: passing to a subgraph never creates a new cycle.

  \item \textbf{Could $H-a$ in Corollary~5.10 be disconnected, invalidating the contradiction?} No. Connectivity is irrelevant for the contradiction. A disconnected graph with minimum degree at least $3$ and no power-of-two cycle is still a counterexample, since each component has minimum degree at least $3$ and every cycle of every component is also a cycle of the whole graph.

  \item \textbf{Is the $2/3$ inequality sharp?} The proof shows exactly where equality can occur: simultaneously $e_H(A,B)=2|A|=4|B|$ and $2e_H(A)=|A|$. Corollary~5.12 verifies that this is equivalent to the rigid configuration ``all vertices of $B$ have degree $4$ and $H[A]$ is a perfect matching.'' No hidden assumptions are used.

  \item \textbf{Is Lemma~5.13 stronger than mere connectivity of $H$?} Yes. Connectivity only guarantees that each component of $H[A]$ has some neighbor in $B$. Lemma~5.13 proves that no component of $H[A]$ can be sealed off by $B$-vertices that connect only to that component; otherwise one gets a proper induced subgraph with minimum degree at least $3$.

  \item \textbf{Interaction with Round~3.} Theorem~5.11 strictly strengthens Round~3 Lemma~5.3: the lower bound on degree-$3$ vertices improves from $4|V(H)|/7$ to $2|V(H)|/3$. The earlier bound remains true but is now superseded.
\end{itemize}

\subsection*{7) FINAL}
\textbf{UNRESOLVED (BUT STRICTLY ADVANCED).} The full Erd\H{o}s--Gy\'arf\'as conjecture remains open in general, but Round~4 makes a genuine structural advance beyond Round~3:
\begin{itemize}
  \item it upgrades edge-minimality to full subgraph-minimality (Lemma~5.9);
  \item it proves that every degree-$3$ vertex has a degree-$3$ neighbor (Corollary~5.10);
  \item it strengthens the Round~3 bound from $4/7$ to $2/3$ on the proportion of degree-$3$ vertices (Theorem~5.11);
  \item it characterizes the equality case $|A|=2|B|$ (Corollary~5.12);
  \item it proves that no connected component of $H[A]$ can be closed under all incident $B$-neighbors unless $H[A]$ itself is connected (Lemma~5.13).
\end{itemize}
The sharpest obstruction now isolated is this: any minimal counterexample must be a connected graph of diameter at least $3$, containing an induced $P_{13}$, in which at least two-thirds of all vertices have degree exactly $3$, every such vertex has another degree-$3$ neighbor, and the extremal equality case forces the degree-$3$ vertices to form a perfect matching while all remaining vertices have degree exactly $4$.

\subsection*{8) COMPLETION ESTIMATE (MANDATORY)}
COMPLETION: 68\%.

\subsection*{9) REFERENCES}
\begin{thebibliography}{99}

\bibitem{Erdos1997}
P.~Erd\H{o}s,
\newblock Some old and new problems in various branches of combinatorics,
\newblock \emph{Discrete Mathematics} \textbf{165--166} (1997), 227--231.

\bibitem{Carr2025}
A.~Carr,
\newblock Cycles of length 4 or 8 in graphs with diameter 2 and minimum degree at least 3,
\newblock arXiv:2508.19302, 2025.

\bibitem{HegdeSandeepShashank2025}
A.~S. Hegde, R.~B. Sandeep, and P.~Shashank,
\newblock Erd\H{o}s--Gy\'arf\'as conjecture on graphs without long induced paths,
\newblock arXiv:2410.22842v2, 2025.

\bibitem{HuShen2024}
Z.~Hu and C.~Shen,
\newblock The Erd\H{o}s--Gy\'arf\'as conjecture holds for $P_{10}$-free graphs,
\newblock \emph{Discrete Mathematics} \textbf{347} (2024), 114175.

\bibitem{NowEtAl2014}
P.~Salehi Nowbandegani, H.~Esfandiari, M.~H. Shirdareh Haghighi, and K.~Bibak,
\newblock On the Erd\H{o}s--Gy\'arf\'as conjecture in claw-free graphs,
\newblock \emph{Discussiones Mathematicae Graph Theory} \textbf{34} (2014), 635--640.

\end{thebibliography}
